\clearpage
\renewcommand{\sectioncolor}{alert}
\section{The evidence trap --- and why your notebook opens with it}
\markboth{The evidence trap}{}

Cells 1--2 of the notebook are logic. But §0, the very first markdown cell, exists to break one
habit before it forms.

\begin{warnbox}[title={The habit: mistaking verification for proof}]
Your companion notebook prints
\begin{cellbox}
ns=3:  Tr(T\^{}ns) == lam\_+\^{}ns + lam\_-\^{}ns  ->  True\\
ns=5:  ...  True\\
ns=8:  ...  True
\end{cellbox}
\vspace{4pt}
Three \bk{True}s feel conclusive. They are three facts about three numbers. The theorem claims
something about \textbf{every} $N$, of which there are infinitely many.
\end{warnbox}

\begin{bookbox}[title={Devlin's killer example, §3.1 --- read this before you ever trust a pattern}]
\begin{quote}\itshape
``Proving a mathematical statement is much more than gathering evidence in its favor\ldots\ the
Goldbach Conjecture\ldots\ It is possible to run computer programs to check the statement for many
specific even numbers, and to date (July 2012) it has been verified for all numbers up to and beyond
$1.6\times10^{18}$ (1.6 quintillion). Most mathematicians believe it to be true. \textbf{But it has
not yet been proved.}''
\end{quote}
\end{bookbox}

\begin{center}
\begin{tikzpicture}[font=\footnotesize]
 \draw[slate!55,line width=1.2pt,-{Stealth[length=5pt]}] (0,0)--(15.2,0);
 \foreach \x/\l in {0.6/{$10^{0}$},3.2/{$10^{6}$},5.8/{$10^{12}$},8.4/{$1.6\times10^{18}$}}{
  \draw[slate!55] (\x,-0.12)--(\x,0.12); \node[font=\scriptsize,below=3pt] at (\x,-0.1) {\l};}
 \fill[bio,opacity=0.22,rounded corners=2pt] (0.3,-0.3) rectangle (8.4,0.3);
 \node[font=\scriptsize\bfseries,text=bio] at (4.3,0.62) {verified by computer};
 \fill[alert,opacity=0.14,rounded corners=2pt] (8.4,-0.3) rectangle (15.0,0.3);
 \node[font=\scriptsize\bfseries,text=alert] at (11.7,0.62) {unknown --- infinitely many cases remain};
 \draw[alert,line width=1.4pt,dashed] (8.4,-0.75)--(8.4,0.95);
 \node[font=\scriptsize,text=slate,anchor=north,text width=13cm,align=center] at (7.6,-1.05)
  {Goldbach's conjecture. \textbf{No finite amount of green closes the red.}
   The same picture applies to \bk{ns=3,5,8}.};
\end{tikzpicture}
\end{center}
\vspace{10pt}

\begin{bookbox}[title={Diedrichs \& Lovett name the error, §2.3.6}]
\begin{quote}\itshape
``\textbf{Try a Few Examples:} Sometimes, running examples with the objects in question allows us to
discern patterns\ldots\ The patterns generally do not constitute a proof, so simply claiming `This
pattern continues' is a form of \textbf{faulty generalization}.''
\end{quote}
So examples are \emph{encouraged} --- as a way to find the proof. They are just never the proof.
\end{bookbox}

\begin{dobox}[title={The one sentence to carry}]
\begin{center}\itshape\large
Checking is a search. Proving is an argument.\\
A search over an infinite set never terminates.
\end{center}
\end{dobox}
